\documentclass[10.5pt,a4paper]{article}

% ---------- Packages ----------
\usepackage[utf8]{inputenc}
\usepackage[T1]{fontenc}
\usepackage[french]{babel}
\usepackage{XCharter}                % elegant serif body font
\usepackage[a4paper,margin=0.85cm]{geometry}
\usepackage{xcolor}
\usepackage{titlesec}
\usepackage{enumitem}
\usepackage{ragged2e}
\usepackage{paracol}
\usepackage[none]{hyphenat}
\setlength{\emergencystretch}{3em}

% ---------- Colors ----------
\definecolor{accent}{RGB}{30,64,149}
\definecolor{graytext}{RGB}{90,90,90}

% ---------- Section titles (main column) ----------
\titleformat{\section}
  {\color{accent}\bfseries\Large}
  {}{0em}{}
\titlespacing{\section}{0pt}{4pt}{2pt}
\newcommand{\sectionrule}{\vspace{-6pt}{\color{accent}\hrule height 1pt}\vspace{3pt}}

% ---------- Sidebar section titles ----------
\newcommand{\subsectitle}[1]{%
  {\color{accent}\bfseries\large #1}\\[-2pt]
  {\color{accent}\hrule height 0.8pt}\vspace{4pt}
}

\setlength{\parindent}{0pt}
\pagestyle{empty}

% ---------- Custom commands ----------

% Experience / generic entry: #1 position #2 org #3 location #4 period #5 body
\newcommand{\cvexp}[5]{%
  \noindent\textbf{#1} \hfill {\itshape\color{graytext}\small #4}\\
  {\itshape #2 --- #3}\\[-6pt]
  #5
  \vspace{9pt}
}

% Education entry: #1 school #2 area/degree #3 location #4 period #5 description
\newcommand{\cvedu}[5]{%
  \noindent\textbf{#1} \hfill {\itshape\color{graytext}\small #4}\\
  {\itshape #2, #3}\\[1pt]
  {\small #5}
  \vspace{2pt}
}

% Education entry without description
\newcommand{\cveduplain}[4]{%
  \noindent\textbf{#1} \hfill {\itshape\color{graytext}\small #4}\\
  {\itshape #2, #3}
  \vspace{2pt}
}

% Project entry: #1 name #2 period #3 description
\newcommand{\cvproject}[3]{%
  \noindent\textbf{#1} \hfill {\itshape\color{graytext}\small #2}\\[1pt]
  {\small #3}
  \vspace{3pt}
}

% Sidebar simple line entry: #1 title (bold) #2 detail
\newcommand{\sidebarline}[2]{%
  \noindent\textbf{\small #1}\\
  {\small #2}\\[3pt]
}

% Sidebar dated entry (volunteer, with description)
\newcommand{\sidebardated}[3]{%
  \textbf{\small #1} {\itshape\color{graytext}\footnotesize (#2)}\\
  {\small #3}\\[3pt]
}

% Sidebar dated entry without description (certifications)
\newcommand{\sidebarcert}[2]{%
  \noindent\textbf{\small #1} {\itshape\color{graytext}\footnotesize (#2)}\\[2pt]
}

\begin{document}

% ================= HEADER =================
\begin{center}
  {\fontsize{26}{30}\selectfont\bfseries Amélie Yang}\\[4pt]
  {\color{accent}\itshape\large Étudiante ingénieure en IA Générative \& Développement Backend -- IMT Mines Albi}\\[6pt]
  {\small
    Grenoble, France
    \,\textbullet\,
    {\color{accent}amelie.yang@mines-albi.fr}
    \,\textbullet\,
    +33~7~67~95~35~99
    \,\textbullet\,
    {\color{accent}linkedin.com/in/amelieyang}
  }
\end{center}

\vspace{3pt}
{\color{accent}\hrule height 1.2pt}
\vspace{4pt}

% ================= TWO COLUMN BODY =================
\columnratio{0.65}
\setlength{\columnsep}{1.5em}
\begin{paracol}{2}

% ---------- RÉSUMÉ ----------
\section*{Résumé}
\sectionrule
\justifying
Ingénieure diplômée de l'IMT Mines Albi (Data \& Systèmes d'Information), avec une expérience concrète de développement de services d'IA générative en production : APIs backend (FastAPI), pipelines RAG, orchestration de LLM et agents IA. Fort attrait pour l'intérêt général et le service public. En recherche d'un poste d'ingénieure IA à partir de septembre 2026.

% ---------- EXPÉRIENCE ----------
\section*{Expérience}
\sectionrule

\cvexp{Stagiaire Ingénieure IA Générative}{ALTEN Lab}{Grenoble, France}{Février 2026 -- Août 2026 (6 mois)}{%
  \begin{itemize}[leftmargin=12pt, itemsep=2pt, topsep=0pt, partopsep=0pt, parsep=0pt]
    \item Conception et développement d'un service RAG multimodal \emph{as-a-service} pour la recherche documentaire industrielle : API backend FastAPI conteneurisée (Docker), pensée comme abstraction réutilisable et adaptable à différents métiers.
    \item Orchestration de LLM, deux moteurs RAG complémentaires (frameworkless, LangChain/\allowbreak LangGraph) ; ingestion (chunking, embeddings), recherche hybride BM25/\allowbreak sémantique (fusion RRF, PostgreSQL : pgvector + ParadeDB).
    \item Qualité du code : benchmark d'outils d'extraction (Docling, PyMuPDF, OCR, VLM), tests, évaluation sur benchmark académique (RAGAS, recall@k, MRR, bootstrap), CI/CD, déploiement sur Open WebUI.
  \end{itemize}
}

\cvexp{Stagiaire Assistante Ingénieure R\&D}{Yooz}{Aimargues, France}{Avril 2025 -- Août 2025 (4 mois)}{%
  \begin{itemize}[leftmargin=12pt, itemsep=2pt, topsep=0pt, partopsep=0pt, parsep=0pt]
    \item Développement d'agents IA « comptable » connectés à un RAG indexé sur le Plan Comptable Général, pensés pour faciliter leur intégration par des clients tiers.
    \item Application web (Dash) de visualisation et d'expérimentation de solutions IA (recherche sémantique, RAG, LLM), pensée pour être prise en main par des utilisateurs métier non techniques.
    \item Clustering et réduction de dimension (PCA, t-SNE, UMAP) sur bases vectorielles (Milvus, ChromaDB) pour diagnostiquer et améliorer la qualité de la recherche sémantique en production.
  \end{itemize}
}

\cvexp{Stagiaire Opératrice}{Laboratoires Pierre Fabre}{Soual, France}{Février 2024 -- Mars 2024 (2 mois)}{%
  \begin{itemize}[leftmargin=12pt, itemsep=2pt, topsep=0pt, partopsep=0pt, parsep=0pt]
    \item Renforcement des compétences opérationnelles en gérant la ligne de conditionnement, surveillance, tri et alimentation des processus.
  \end{itemize}
}

% ---------- FORMATION ----------
\section*{Formation}
\sectionrule

\cvedu{The Catholic University of Korea}{Spécialité IA -- Étudiante en échange}{Séoul, Corée du Sud}{Septembre 2025 -- Décembre 2025}{Computer Vision, Data Mining, Machine Learning.}

\cvedu{IMT Mines Albi}{Ingénieur Généraliste}{Albi, France}{2023 -- 2026}{Optimisation, Base de données, Calcul numérique, Big Data, Machine Learning.}

\cveduplain{Lycée Joffre}{Classe préparatoire MPSI-MP}{Montpellier, France}{2021 -- 2023}

% ---------- PROJETS ----------
\section*{Projets}
\sectionrule

\cvproject{CNN pour la reconnaissance de chiffres manuscrits}{2024}{Développement d'un réseau de neurones convolutif (CNN) pour la classification d'images (Dataset MNIST) avec PyTorch : couches de convolution/pooling, fonctions d'activation, rétropropagation, optimisation (Adam) et régularisation (dropout).}

\cvproject{Prévision en temps réel du temps d'attente à quai de la SNCF}{2025}{Nettoyage, traitement et agrégation des données de flux. Entraînement de modèles de régression et de séries temporelles en Python.}

\switchcolumn

\vspace{2pt}

% ---------- COMPÉTENCES ----------
\subsectitle{Compétences}

\sidebarline{Langages}{Python (FastAPI, Pandas, NumPy, Scikit-learn, PyTorch), SQL, C++}
\sidebarline{IA générative / LLM}{RAG (as-a-service), orchestration de LLM, agents IA, LangChain, LangGraph, LlamaIndex, embeddings}
\sidebarline{Backend \& Infrastructure}{FastAPI, architectures d'intégration, Docker, CI/CD, Git}
\sidebarline{Bases de données}{SQL, BM25 (ParadeDB), bases vectorielles (pgvector, Milvus, ChromaDB)}
\sidebarline{Machine Learning}{Classification, clustering, réduction de dimension, CNN, séries temporelles}

% ---------- LANGUES ----------
\subsectitle{Langues}
\noindent{\small Anglais (TOEIC 965 pts), Espagnol, Japonais}

\vspace{5pt}

% ---------- ASSOCIATIF ----------
\subsectitle{Associatif}

\sidebardated{Secrétaire Générale -- Bureau des Sports IMT Mines Albi}{2024 -- 2025}{Organisation de plannings et d'événements. Rédaction de dossiers, comptes-rendus.}

\sidebardated{Les Cordées de la Réussite}{2023 -- 2024}{Encadrement et accompagnement d'élèves dans leur orientation.}

% ---------- CERTIFICATIONS ----------
\subsectitle{Certifications}

\sidebarcert{MOOC -- Prévention des risques professionnels (INRS)}{2023}
\sidebarcert{MOOC -- Propriété Intellectuelle (INPI)}{2024}

\vspace{4pt}

% ---------- INTÉRÊTS ----------
\subsectitle{Intérêts}
\noindent{\small Volley (2 ans en compétition)}\\
{\small Musique (piano, guitare)}\\
{\small Jeux-vidéos}

\end{paracol}

\end{document}
